\chapter{Einleitung}

Die fortschreitende Digitalisierung hat in den vergangenen Jahren zur Entstehung
komplexer, datengetriebener Plattformökosysteme geführt, die mittlerweile nahezu alle
Bereiche des gesellschaftlichen, wirtschaftlichen und politischen Lebens durchdringen.
Soziale Netzwerke, Suchmaschinen und digitale Plattformen wie Google, TikTok, Meta oder
X setzen zunehmend \ac{KI}-gestützte Systeme ein, um Nutzerverhalten zu erfassen, zu
analysieren und gezielt zu beeinflussen. Durch algorithmische Steuerung (Algorithmic
Steering) und Engagement-Optimierung werden Inhalte personalisiert, Aufmerksamkeit
kanalisiert und Interaktionsmuster systematisch gelenkt.

Diese Entwicklung geht weit über reine Informationsbereitstellung hinaus: Moderne
\ac{KI}-gestützte Personalisierungssysteme (AI-driven Personalization Systems) erzeugen
individualisierte Informationsumgebungen, die Wahrnehmung, Meinungsbildung und
Entscheidungsprozesse der Nutzer nachhaltig prägen können. Die daraus resultierende
Verhaltensbeeinflussung (Behavioural Influence) wirft grundlegende Fragen hinsichtlich
individueller Autonomie, informationeller Selbstbestimmung und demokratischer
Willensbildung auf.

Gleichzeitig zeigen aktuelle Debatten um Desinformationsdynamiken (Misinformation
Dynamics) und soziale Fragmentierung (Social Fragmentation), dass algorithmisch gesteuerte
Plattformen zur Polarisierung öffentlicher Diskurse und zur Erosion gemeinsamer
Wissensgrundlagen beitragen können. Im Kontext des Überwachungskapitalismus (Surveillance
Capitalism) werden diese Mechanismen zudem für ökonomische Verwertungslogiken
instrumentalisiert, während staatliche Akteure ähnliche Technologien zur politischen
Manipulation und sozialen Kontrolle einsetzen.

Besonders problematisch erscheint dabei, dass viele dieser Systeme bislang kaum reguliert
sind (Barely Regulated). Trotz ihrer weitreichenden gesellschaftlichen Auswirkungen
operieren algorithmische Monitoring- und Bewertungssysteme häufig in regulatorischen
Grauzonen, in denen weder ausreichende Transparenzpflichten noch wirksame
Kontrollmechanismen existieren.

Vor diesem Hintergrund untersucht die vorliegende Arbeit, wie \ac{KI}-gestützte
Monitoring-Systeme als digitale Informationssysteme funktionieren und inwiefern sie
datenbasierte Analyse und Beeinflussung menschlichen Verhaltens ermöglichen.


% ============================================================
\section{Problemstellung}
% ============================================================

Die zunehmende Durchdringung digitaler Lebenswelten durch algorithmische Systeme hat eine
Situation geschaffen, in der Verhaltensbeeinflussung nicht mehr als Ausnahme, sondern als
strukturelles Merkmal moderner Plattformökonomien betrachtet werden muss. Digitale
Informationssysteme und \ac{KI}-gestützte Plattformen erfassen kontinuierlich
nutzerbezogene Daten und setzen diese mittels \ac{ML}-basierter Analyseverfahren ein, um
personalisierte Inhalte zu generieren, Nutzerverhalten zu prognostizieren und
Interaktionen gezielt zu steuern.

Die dabei eingesetzten Mechanismen der Engagement-Optimierung zielen primär auf die
Maximierung von Verweildauer, Interaktionsraten und datenbasierter Monetarisierung ab.
Empfehlungssysteme, Ranking-Algorithmen und personalisierte Werbesysteme bilden dabei die
technologische Grundlage einer umfassenden algorithmischen Steuerung, die weit über
neutrale Informationsvermittlung hinausgeht.

Parallel dazu gewinnen staatlich geprägte Überwachungs- und Bewertungssysteme zunehmend
an Bedeutung. Insbesondere im Kontext des chinesischen \ac{SCS} wird deutlich, wie
datenbasierte Monitoring-Infrastrukturen mit Mechanismen sozialer Bewertung und
Verhaltensregulierung verknüpft werden können. Die Konvergenz kommerzieller und
staatlicher Überwachungslogiken verdeutlicht die gesellschaftspolitische Tragweite
algorithmischer Monitoring-Systeme.

Aus wissenschaftlicher Perspektive ergeben sich hieraus mehrere zentrale Problemfelder:
Die Intransparenz proprietärer Algorithmen erschwert eine unabhängige Bewertung ihrer
Wirkungsmechanismen. Desinformationsdynamiken und soziale Fragmentierung werden durch
algorithmische Verstärkungseffekte begünstigt. Die bestehenden regulatorischen Rahmen
erweisen sich als unzureichend, um den Herausforderungen algorithmischer
Verhaltensbeeinflussung angemessen zu begegnen.

Die vorliegende Arbeit adressiert diese Problemstellung durch eine systematische
Untersuchung der Funktionsweise, Wirkungsmechanismen und gesellschaftlichen Implikationen
\ac{KI}-gestützter Monitoring-Systeme.


% ============================================================
\section{Zielsetzung und Forschungsfrage}
% ============================================================

Ziel der vorliegenden Arbeit ist die systematische Untersuchung \ac{KI}-gestützter
Monitoring-Systeme als digitale Informationssysteme sowie die Analyse ihrer Möglichkeiten
zur datenbasierten Erfassung, Bewertung und Beeinflussung menschlichen Verhaltens.

Im Mittelpunkt stehen dabei folgende Untersuchungsschwerpunkte: die technologischen
Grundlagen algorithmischer Steuerung und Engagement-Optimierung, die Mechanismen der
Verhaltensbeeinflussung durch \ac{KI}-gestützte Personalisierungssysteme, die Rolle von
Desinformationsdynamiken und sozialer Fragmentierung im Kontext algorithmischer
Plattformlogiken, die ökonomischen Verwertungsstrukturen des Überwachungskapitalismus
sowie die Möglichkeiten politischer Manipulation durch datenbasierte Monitoring-Systeme.

Die Arbeit verfolgt dabei einen literaturbasierten Forschungsansatz, der theoretische
Grundlagen mit einer kritischen Diskussion verbindet. Die Untersuchung gliedert sich in
theoretische Kapitel zur technologischen und konzeptionellen Fundierung, eine
systematische Literaturanalyse relevanter Forschungsbeiträge sowie Diskussionsabschnitte,
in denen die gewonnenen Erkenntnisse in kohärente Argumentationsketten überführt werden.

Darüber hinaus werden mögliche Hypothesen zur Wirkungsweise algorithmischer
Monitoring-Systeme formuliert, die als Grundlage für weiterführende empirische Forschung
dienen können.

Die zentrale Forschungsfrage der Arbeit lautet:

\enquote{Wie funktionieren \ac{KI}-gestützte Monitoring-Systeme als digitale
Informationssysteme und inwiefern ermöglichen sie datenbasierte Analyse und Beeinflussung
menschlichen Verhaltens?}

Zur Beantwortung der Forschungsfrage erfolgt eine qualitative Literaturrecherche und
theoriegeleitete Analyse wissenschaftlicher Fachliteratur aus den Bereichen künstliche
Intelligenz, digitale Informationssysteme, Plattformökonomie, algorithmische Überwachung
sowie Verhaltens- und Kommunikationsforschung.


% ============================================================
\section{Aufbau der Arbeit}
% ============================================================

Die vorliegende Bachelorarbeit ist in mehrere thematische Bestandteile gegliedert, die
aufeinander aufbauend die Forschungsfrage systematisch bearbeiten.

Kapitel~2 behandelt die theoretischen und technologischen Grundlagen der Arbeit. Zunächst
werden zentrale Begriffe und Konzepte aus den Bereichen \ac{KI}, \ac{ML} und digitale
Informationssysteme erläutert. Anschließend werden Monitoring- und Bewertungssysteme
sowie relevante Verfahren der Datenerfassung, algorithmischen Analyse und Klassifikation
dargestellt. Darüber hinaus wird das methodische Vorgehen der Arbeit beschrieben,
einschließlich der Literaturrecherche, der Auswahlkriterien der Quellen sowie
methodischer Grenzen.

Kapitel~3 untersucht \ac{KI}-gestützte Monitoring-Systeme in digitalen
Informationsumgebungen. Der Schwerpunkt liegt auf der Datenerfassung und digitalen
Verhaltensanalyse innerhalb moderner Plattformökosysteme. Zudem werden algorithmische
Steuerungs- und Bewertungsmechanismen wie Personalisierung, Ranking-Algorithmen,
Reputations- und Scoring-Systeme sowie Engagement-Optimierung und Verhaltensprognosen
analysiert. Abschließend erfolgt ein Vergleich unterschiedlicher Systemlogiken,
insbesondere zwischen plattformzentrierten und staatlich integrierten Monitoring-Systemen.

Kapitel~4 beschäftigt sich mit den sozioökonomischen und gesellschaftlichen Implikationen
algorithmischer Monitoring-Systeme. Dabei werden Auswirkungen auf digitale
Informationsumgebungen, Verhaltensdynamiken und Personalisierungseffekte, ökonomische
Interessen und Plattformlogiken sowie politische und gesellschaftliche
Einflussmechanismen betrachtet. Darüber hinaus werden Governance- und Kontrollstrukturen
diskutiert.

Kapitel~5 bietet eine kritische Betrachtung der Arbeit, einschließlich der Grenzen der
technologischen Analyse, der Literaturbasis sowie der Aussagekraft und methodischen
Limitationen. Zudem werden normative und empirische Perspektiven reflektiert.

Kapitel~6 fasst die zentralen Erkenntnisse zusammen, beantwortet die Forschungsfrage und
gibt einen Ausblick auf zukünftige Entwicklungen im Bereich \ac{KI}-gestützter
Monitoring-Systeme und digitaler Informationssysteme.
